\section{Introducci\'on}

En este trabajo se busca investigar e implementar un m\'etodo de busqueda
vectorial usada en las bases de datos vectoriales. 
Para esto, se tom\'o como referencia la librer\'ia popular de busqueda vectorial 
FAISS \cite{faiss_gh}, especialmente un m\'etodo de b\'usqueda 
de vecino m\'as cercano (NNS, Nearest Neighbour Search) con cuantificacio\'on
de productos (PQ) \cite{pqnns} que citan en sus
\href{https://faiss.ai/index.html}{docs}.

El motivo de este trabajo es comprender la escencia en la b\'usqueda 
sem\'antica, una herramienta que permite a la computadora derivar significado
de palabras y oraciones, junto con la optimizaci\'on espacial que aporta el 
algoritmo implementado.

Ademas, la implementaci\'on esta hecha en C, en un solo documento
\verb|jumpsuit.h|\footnote{El nombre 'jumpsuit' y 'orange-jumpsuit' son por la
vestimenta del personaje 'Vector' en la pelicula 'Mi Villano Favorito'},
con minimas dependencias, para que un usuario de la librer\'ia pueda copiar 
y pegar el documento en su proyecto sin mayor problema. 
Tambi\'en se provee un wrapper de Python para que se pueda usar como reemplazo 
b\'asico del wrapper de python que provee FAISS. 
Dados estos objetivos de experiencia de usuario, y el tiempo de desarrollo, 
la librer\'ia es limitada, en aspectos de funcionalidad y eficiencia (memoria y 
velocidad).


